\documentclass{situationist-preprint}

\title{Two Epistemic Statuses}
\msid{EPI-06-a}
\msfamily{EPI-06}
\msdiscipline{Epistemology}
\mstype{Principal manuscript}
\mskeywords{internalism; externalism; justification; knowledge; new evil demon; clairvoyance; mentalism; safety; entitlement}

\begin{document}
\maketitle

\begin{abstract}
Is epistemic justification determined by factors internal to the subject's mind, or can it depend on facts, such as the reliability of her faculties, to which she has no access? Each side has a case that the other cannot accommodate: the deceived victim of a demon seems as justified as we are, which favours internalism, and a belief can be reliably formed and yet irresponsibly held, which tells against identifying justification with reliability. We argue that the dispute has been conducted over two statuses under one name. Justification answers the question whether a subject has believed as she should, given what she has to go on; knowledge answers the question whether her belief can be relied on as a source of truth. We show that a status answering to the first concern must be shared by internal duplicates, and that a status answering to the second cannot be. Hence internalism holds for justification and externalism for knowledge. The thesis fits the analysis of knowledge as safe, justified, true belief, explains why justification is not context-relative, and leaves room for the externalist standing of rules of inference on which the justification of induction depends.
\end{abstract}

\section{The problem}

Internalism about justification holds that whether a subject is justified in believing $p$ is fixed by factors internal to her: on access internalism, factors to which she has reflective access \citep{BonJour1985}; on mentalism, her mental states \citep{ConeeFeldman2001}. Externalism denies this. On process reliabilism, a belief is justified if it is produced by a reliable belief-forming process \citep{Goldman1979}; on proper functionalism and its relatives, if it is produced by faculties functioning properly in an appropriate environment \citep{Bergmann2006}.

Two cases have dominated the debate. The first is the \emph{new evil demon}. A subject whose experiences are systematically produced by a deceiving demon, and are exactly like ours, forms beliefs just as we do. Her beliefs are formed by unreliable processes, but she seems to be as justified as we are \citep{LehrerCohen1983,Cohen1984}. The second is \emph{clairvoyance}. \citet{BonJour1980} described Norman, who has a reliable clairvoyant power of which he has no evidence, and who one day finds himself believing that the President is in New York. Norman's belief is reliably produced, but, holding it without any reason he could give and against the background of his ignorance of his power, he seems to be epistemically irresponsible, and so unjustified.

Externalists have answered the first case by distinguishing kinds of justification. \citet{Goldman1988} distinguished strong justification, which requires reliability, from weak justification, which is blameless belief; \citet{Engel1992} distinguished personal justification, an evaluation of the subject, from doxastic justification, an evaluation of the belief; \citet{Sosa1991} distinguished the aptness of a belief from its being justified from the subject's perspective. Internalists have answered by holding that reliability is relevant to knowledge but not to justification \citep{ConeeFeldman2004}, and \citet{Alston1988} defended a hybrid, on which the ground of a justified belief must be accessible but its adequacy need not be.

We argue for a thesis that makes these distinctions principled. There are two epistemic statuses: justification, which is internalist, and knowledge, which is externalist. They are not rival analyses of one status, and the cases that seem to support one side against the other bear on different statuses. The thesis is stated as a supervenience claim and derived from an account of the concerns the two statuses answer to. We then show that it is the account of justification presupposed by the analysis of knowledge in EPI-01 and used in EPI-02, EPI-13, EPI-19 and EPI-21, and we explain how it is consistent with the reliabilist account of the justification of induction in EPI-03.

\section{Framework}

\begin{definition}[Internal duplicates]\label{d:dup}
Two subjects, in the same or different worlds, are \emph{internal duplicates} at a time if they are alike in all their mental states at that time and in the history of those states: in how things seem to them, perceptually, mnemonically and intellectually; in what they believe and on what bases; and in what they have attended to and inferred.
\end{definition}

The definition is mentalist rather than accessibilist. It does not require that the subject be able to tell, by reflection, which mental states she is in. Including the history of mental states allows the internal facts to fix whether a belief was formed by inference from other beliefs, and so whether it is based on what supports it; a belief's basis, as the subject has it, is internal.

\begin{definition}[Internalism and externalism]\label{d:int}
An epistemic status $S$ is \emph{internalist} if, necessarily, internal duplicates have $S$ with respect to the same propositions and to the same degree. It is \emph{externalist} otherwise.
\end{definition}

Our thesis concerns two statuses, justification and knowledge, and it derives their character from what each is for. The method is that of \citet{Craig1990}, who asked what the concept of knowledge does for creatures like us, and applied it to the concept of knowledge alone; we apply it to both.

\begin{principle}[The concern of justification]\label{p:jconcern}
To ask whether a subject is justified in believing $p$ is to ask whether, in believing $p$, she has responded as she should to what she has to go on.
\end{principle}

\begin{principle}[The concern of knowledge]\label{p:kconcern}
To ask whether a subject knows $p$ is to ask whether her belief that $p$ can be relied on, by herself and by others, as a source of the truth about whether $p$.
\end{principle}

\ref{p:jconcern} places justification among the evaluations of a subject's conduct as a believer: it is what we assess when we ask whether she reasoned well, whether she ignored evidence, whether she was credulous or careful. Such evaluations take the subject's situation as given and ask how she handled it. \ref{p:kconcern} places knowledge among the evaluations of a subject as an informant. We need to identify those whose beliefs we can rely on, and we need to identify, in our own case, which of our beliefs can be relied on in reasoning and action; the concept of knowledge marks these beliefs \citep{Craig1990}. Its concern is not with the subject's conduct but with the relation between her belief and the fact.

\begin{assumption}[What one has to go on]\label{a:goon}
What a subject has to go on, in forming and holding beliefs, is exhausted by her mental states and their history.
\end{assumption}

Assumption~\ref{a:goon} is not the access requirement. A subject's mental states include states she cannot introspectively discriminate from others \citep[ch.~4]{Williamson2000}, and she has them to go on whether or not she can tell that she has them. What the assumption excludes is that the reliability of her faculties, or the layout of the world beyond her experience of it, is among the things she has to go on. These may explain why her states are as they are; they are not states of hers that she can respond to.

\begin{assumption}[Divergence]\label{a:div}
There are internal duplicates whose beliefs that $p$ differ in whether they can be relied on as a source of the truth about whether $p$.
\end{assumption}

Assumption~\ref{a:div} is illustrated by every case in which a subject and her deceived duplicate believe the same thing on the same basis: the first is reliable on the question, and the second, whose belief is true by luck if it is true at all, is not.

\section{The argument}

\begin{lemma}\label{l:j}
Justification is internalist.
\end{lemma}

\begin{proof}
Let $A$ and $B$ be internal duplicates. By Assumption~\ref{a:goon}, they have the same things to go on. By Definition~\ref{d:dup}, they respond to them in the same way: they hold the same beliefs on the same bases, by the same inferences, having attended to the same considerations. Whether one has responded as one should to what one has to go on depends only on what one has to go on and on how one has responded to it. By \ref{p:jconcern}, then, $A$ is justified in believing $p$, to a given degree, if and only if $B$ is.
\end{proof}

The proof uses one substantive premise beyond the two principles: that the question how well a subject responded to her situation depends only on her situation and her response. This is a general feature of evaluations of conduct. A navigator who plots a correct course from charts that are, unknown to him and undetectably, wrong has navigated as well as one who plots the same course from correct charts; one of them reaches harbour, but the evaluation of the navigation is the same. Outcomes and the reliability of instruments bear on whether we should rely on a navigator's results, not on whether he navigated well.

\begin{lemma}\label{l:k}
Knowledge is externalist.
\end{lemma}

\begin{proof}
By Assumption~\ref{a:div}, there are internal duplicates whose beliefs that $p$ differ in whether they can be relied on as a source of the truth about $p$. By \ref{p:kconcern}, knowledge marks exactly the beliefs that can be so relied on. So one duplicate's belief is knowledge and the other's is not, and by Definition~\ref{d:int} knowledge is externalist.
\end{proof}

\begin{theorem}[Two statuses]\label{t:main}
Justification is internalist and knowledge is externalist. In particular, they are distinct statuses, and there are justified beliefs that are not knowledge for reasons that no internal difference explains.
\end{theorem}

\begin{proof}
The first sentence combines Lemmas~\ref{l:j} and~\ref{l:k}. For the second, take the duplicates of Assumption~\ref{a:div}, and suppose the reliable one is justified and knows. By Lemma~\ref{l:j}, the other is equally justified; by Lemma~\ref{l:k}, she does not know. Her justified belief fails to be knowledge, and nothing internal to her explains the failure, since she is internally like the knower.
\end{proof}

\begin{corollary}[The new evil demon]\label{c:demon}
The demon's victim, and the envatted brain, are justified in their perceptual beliefs to the same degree as their unvictimized duplicates, and do not know.
\end{corollary}

\begin{proof}
They are internal duplicates of ordinary perceivers. By Lemma~\ref{l:j}, they are as justified; their beliefs are not reliable sources of truth, so by \ref{p:kconcern} they are not knowledge.
\end{proof}

\begin{corollary}[Context]\label{c:context}
Justification is not relative to the context of attribution.
\end{corollary}

\begin{proof}
A change in the context in which a justification claim is made, such as the attributor's raising of a sceptical possibility, changes nothing in the mental states of the subject to whom justification is attributed. By Lemma~\ref{l:j}, it changes nothing in whether she is justified.
\end{proof}

The proof of Corollary~\ref{c:context} assumes that a context of attribution can change the subject's justification only by changing something about her; a contextualist about \enquote{justified} who denied this would have to hold that the concern of \ref{p:jconcern} varies with the attributor's interests, and we know of no evidence that it does. Knowledge is different. Whether a belief can be relied on as a source of truth depends on which error possibilities the reliance must be proof against, and that depends on the purposes of those who rely on it. EPI-13 develops this into the thesis that \enquote{know} is context-sensitive through the closeness relation on which safety depends, while justification, as Corollary~\ref{c:context} requires, is not.

\section{How the statuses combine}

The two statuses are not independent. A belief that is relied on as a source of truth is held by someone, and a source whose holder believes irresponsibly is not one we can rely on: Norman, if he knew that his clairvoyant beliefs had no credentials he could appeal to, would have to treat them as hunches, and others who knew how he came by them would do the same. Knowledge, on the analysis of EPI-01, is safe, justified, true belief: it combines the internalist status with an externalist condition, safety, which concerns whether the belief could easily have been false on its basis. Theorem~\ref{t:main} explains the shape of that analysis. The justification clause answers to the subject's conduct, the safety clause to her reliability as a source, and Gettier cases are those in which the first is met and the second is not. Norman's belief, reliable but unjustified, fails the first clause; the victim's belief, justified but unreliable, fails the second.

The statuses also combine in the structure of justification itself. On EPI-19, if it seems to a subject that $p$, she has prima facie justification for believing $p$. Seemings are mental states, and the principle is internalist in the sense of Definition~\ref{d:int}. On EPI-21, one's evidence is one's knowledge, so that evidence is externalist: the perceiver has as evidence that she has hands, and her envatted duplicate does not. The two theses are consistent with Lemma~\ref{l:j} because evidence, so understood, is not what fixes the degree of justification. The perceiver's belief that she has hands and her duplicate's belief are justified by the same seeming; the perceiver's further inferences proceed from premises she knows and her duplicate's from premises she justifiably believes, and justification is transmitted by good inference from justified premises whether or not they are known. The duplicates are equally justified and unequally evidenced, as EPI-21 holds.

\paragraph{The standing of rules.} EPI-03 argues that induction is justified, and justified inductively, given a reliabilist principle on which an inference from justified premises by a rule that is reliable in the actual world, without a defeater, yields a justified conclusion. This may seem to conflict with Lemma~\ref{l:j}. It does not, once two questions are distinguished. The first is whether a subject who reasons inductively has responded as she should to what she has to go on. She has: a competent inductive reasoner, whose inferences seem to her to be good and who has no reason to doubt them, is justified in the sense of \ref{p:jconcern}, and so is her duplicate in a world where induction fails. The second is whether the rule she follows has the standing that makes the beliefs it produces fit to be relied on as sources of truth, which is the question Hume raised when he asked whether inductive inference leads from truths to truths. This is a question about the rule's reliability, and it is externalist. It is the second question that EPI-03 answers, and its reliabilist principle is best read as a claim about the epistemic standing of rules, the entitlement that a reliable rule confers on the beliefs it produces and that a belief needs, in addition to justification, if it is to be safe and so knowledge. EPI-03 marks the distinction itself when it says that a subject with the same experiences in a sceptical scenario is not justified \enquote{or not in the same way}: she is justified in the sense of this manuscript, and her rules lack the standing that EPI-03 concerns. Its no-access principle, that a subject need not be justified in believing her rule reliable for its use to justify, is accepted by mentalism, which imposes no higher-order requirement. The result is a division of labour: the justification of an inductive belief is internalist, the standing of the inductive rule is externalist, and inductive knowledge requires both.

\section{Objections}

\begin{objection}[Two statuses, one word]
Ordinary speakers apply \enquote{justified} to Norman's belief with hesitation and to the victim's without, and they do not speak of two statuses. A theory on which the externalist intuitions concern knowledge and the internalist ones justification legislates rather than describes.
\end{objection}

\begin{reply}
The theory does not claim that \enquote{justified} is ambiguous; it claims that the word names an internalist status, and that the externalist pressure on it comes from its role in the analysis of knowledge. The hesitation about Norman is explained by Lemma~\ref{l:j}: a subject who believes without anything to go on, or against what he has to go on, has not responded as he should, and his reliability does not change that. The willingness to call the victim justified is explained by the same lemma. What externalists have called justification, such as Goldman's strong justification or Sosa's aptness, are conditions that bear on whether a belief is fit to be relied on; we do not dispute their importance, only their name.
\end{reply}

\begin{objection}[Truth-connection]
Justification is supposed to be truth-conducive; that is what makes it epistemic rather than prudential or moral. An internalist status can be had in a world where it is no guide to truth, so it is not epistemic \citep{BonJour1985}.
\end{objection}

\begin{reply}
Justification is truth-directed in its aim, not truth-conducive in every world. The subject who responds as she should to what she has to go on is trying to believe what is true, and her responses are evaluated by standards whose point is truth; that is what makes the evaluation epistemic. Whether her responses in fact lead to truth depends on her world, and the status that tracks that dependence is knowledge, through safety. Requiring of justification what only knowledge can deliver is the error that produced the externalist account of justification, and the demon world is the case that exposes it.
\end{reply}

\begin{objection}[Overintellectualization]
Internalism demands that justifying factors be accessible, and children and animals cannot reflect on their grounds. It therefore denies them justification, which is absurd \citep{Bergmann2006}.
\end{objection}

\begin{reply}
The objection applies to access internalism, not to the mentalist internalism of Definition~\ref{d:int}. A child to whom it looks as if there is a dog is justified in believing there is a dog, by EPI-19, without being able to reflect on her seeming or to represent it as a reason. Bergmann's dilemma, that internalism requires either an awareness of grounds that leads to regress or an awareness too weak to answer the internalist's motivation, presupposes that the motivation of internalism is a demand for awareness. On the present view the motivation is \ref{p:jconcern}, which requires only that justification depend on the subject's states and her response to them.
\end{reply}

\begin{objection}[Anti-luminosity]
No non-trivial mental condition is such that one is always in a position to know whether one is in it \citep{Williamson2000}. If internalism is motivated by the thought that one can always tell whether one is justified, it fails, and with it the case for Lemma~\ref{l:j}.
\end{objection}

\begin{reply}
Lemma~\ref{l:j} is a supervenience claim and makes no claim about what one can tell. A subject may be unable to tell whether it seems to her that the patch is red, and so unable to tell whether she is justified in believing it red; she is nonetheless justified or not according to whether it does seem so. Anti-luminosity tells against the claim that justification is always knowable, which we do not make.
\end{reply}

\begin{objection}[Knowledge first]
If knowledge is prior to justification, as \citet{Williamson2000} argues, a belief is justified only if it is knowledge, or only to the extent that it is supported by what one knows. Then the victim is not justified but at best blameless \citep{Sutton2007,Littlejohn2012}.
\end{objection}

\begin{reply}
On that view the victim must be described as blameless but unjustified, and the defender of the view owes an account of the difference that does not reintroduce an internalist status under another name. Blamelessness so described is shared by internal duplicates, is assessed by the subject's response to what she has to go on, and has degrees that track how carefully she reasoned; it is the status of \ref{p:jconcern}. Whether it is called justification or excuse, there are two statuses, one internalist and one externalist, which is our thesis. We keep the traditional name for the first because it is the status that EPI-01 combines with safety and that EPI-05 assigns a tree structure.
\end{reply}

\section{What the result does not establish}

Theorem~\ref{t:main} does not say which mental states justify and how; EPI-19 defends phenomenal conservatism for seemings, and EPI-05 the structure of support. It does not analyse knowledge, which is the task of EPI-01; it shows only that any adequate analysis will be externalist. It takes no position on whether justification is the only internalist epistemic status, and it leaves to EPI-06-b the detailed treatment of the variants of the new evil demon problem. Finally, \ref{p:jconcern} and \ref{p:kconcern} are claims about the concerns the two concepts serve. A theorist who holds that one concept, under the name \enquote{justification}, answers to both will find that it is neither internalist nor externalist, and the cases of the demon and of Norman will remain, for her, a standing conflict of intuitions.

\section{Conclusion}

Justification answers to the question whether a subject has responded as she should to what she has to go on, and so cannot differ between internal duplicates; knowledge answers to the question whether her belief can be relied on as a source of truth, and so must. The demon's victim is justified and does not know; Norman's reliable belief is not justified and so is not knowledge. Knowledge combines the two concerns, as the analysis of knowledge as safe, justified, true belief requires, and the reliability of rules that EPI-03 invokes is part of the externalist standing a belief needs, beyond justification, to be knowledge.

\begin{thebibliography}{99}\small
\bibitem[Alston(1988)]{Alston1988} Alston, W.~P. (1988). An internalist externalism. \emph{Synthese}, 74(3), 265--283.
\bibitem[Bergmann(2006)]{Bergmann2006} Bergmann, M. (2006). \emph{Justification without Awareness: A Defense of Epistemic Externalism}. Oxford University Press.
\bibitem[BonJour(1980)]{BonJour1980} BonJour, L. (1980). Externalist theories of empirical knowledge. \emph{Midwest Studies in Philosophy}, 5, 53--73.
\bibitem[BonJour(1985)]{BonJour1985} BonJour, L. (1985). \emph{The Structure of Empirical Knowledge}. Harvard University Press.
\bibitem[Cohen(1984)]{Cohen1984} Cohen, S. (1984). Justification and truth. \emph{Philosophical Studies}, 46(3), 279--295.
\bibitem[Conee and Feldman(2001)]{ConeeFeldman2001} Conee, E. and Feldman, R. (2001). Internalism defended. \emph{American Philosophical Quarterly}, 38(1), 1--18.
\bibitem[Conee and Feldman(2004)]{ConeeFeldman2004} Conee, E. and Feldman, R. (2004). \emph{Evidentialism: Essays in Epistemology}. Oxford University Press.
\bibitem[Craig(1990)]{Craig1990} Craig, E. (1990). \emph{Knowledge and the State of Nature: An Essay in Conceptual Synthesis}. Oxford University Press.
\bibitem[Engel(1992)]{Engel1992} Engel, M. (1992). Personal and doxastic justification in epistemology. \emph{Philosophical Studies}, 67(2), 133--150.
\bibitem[Goldman(1979)]{Goldman1979} Goldman, A.~I. (1979). What is justified belief? In G.~S. Pappas (Ed.), \emph{Justification and Knowledge} (pp.~1--23). Reidel.
\bibitem[Goldman(1988)]{Goldman1988} Goldman, A.~I. (1988). Strong and weak justification. \emph{Philosophical Perspectives}, 2, 51--69.
\bibitem[Lehrer and Cohen(1983)]{LehrerCohen1983} Lehrer, K. and Cohen, S. (1983). Justification, truth, and coherence. \emph{Synthese}, 55(2), 191--207.
\bibitem[Littlejohn(2012)]{Littlejohn2012} Littlejohn, C. (2012). \emph{Justification and the Truth-Connection}. Cambridge University Press.
\bibitem[Sosa(1991)]{Sosa1991} Sosa, E. (1991). \emph{Knowledge in Perspective: Selected Essays in Epistemology}. Cambridge University Press.
\bibitem[Sutton(2007)]{Sutton2007} Sutton, J. (2007). \emph{Without Justification}. MIT Press.
\bibitem[Williamson(2000)]{Williamson2000} Williamson, T. (2000). \emph{Knowledge and Its Limits}. Oxford University Press.
\end{thebibliography}

\end{document}
